% ৪. নিম্নস্তরের ভাষা
\section{নিম্নস্তরের ভাষা}

\begin{learningbox}
এই টপিক শেষে শিক্ষার্থী low-level language-এর ধরন, সুবিধা ও সীমাবদ্ধতা লিখতে পারবে।
\end{learningbox}

\begin{definitionbox}{নিম্নস্তরের ভাষা}
যে programming language hardware বা machine-এর কাছাকাছি এবং machine dependent, তাকে নিম্নস্তরের ভাষা বলা হয়। Machine language ও assembly language এর অন্তর্ভুক্ত।
\end{definitionbox}

\begin{center}
\fbox{\parbox{0.88\textwidth}{
\centering
\textbf{Tree placeholder}\\[4pt]
Low Level Language $\rightarrow$ Machine Language, Assembly Language
}}
\end{center}

\begin{keypointbox}{মনে রাখো}
\begin{itemize}
\item execution দ্রুত হতে পারে।
\item hardware control বেশি।
\item লেখা, বুঝা ও debug করা কঠিন।
\item machine পরিবর্তন হলে program বদলাতে হতে পারে।
\end{itemize}
\end{keypointbox}
